\documentclass[11pt]{letter}
\usepackage[letterpaper,margin=1in]{geometry}
\usepackage{mathptmx}
\usepackage{microtype}
\signature{Daryl Janzen}
\address{Department of Physics and Engineering Physics\\
University of Saskatchewan\\
116 Science Place\\
Saskatoon, SK S7N 5E2, Canada\\
daryl.janzen@usask.ca}
\date{\today}
\begin{document}
\begin{letter}{The Editors\\
\textit{Cognitive Science}}
\opening{Dear Editors,}

I submit for consideration as a Regular Article the manuscript ``The Assimilate--Receipt--Pivot: A Constitutive Mechanism of Conversation.''

The manuscript proposes that conversation, as distinct from lecture, monologue and parallel speech, is constituted by the exchange of a single move with three components: the listener assimilates the prior contribution, demonstrates the assimilation in a public receipt, and pivots to new content the assimilation made available. Moves are located in a two-axis space: participation, signed so that moves against the conversational contract are distinguished from activity that simply does not engage it, and quality, the degree to which the three components are compressed. The manuscript argues that the components are mutually constitutive, shows that the construct reduces to none of its parent theories, and derives eleven predictions, several of which follow from the integrated construct alone, together with a programme for testing them.

The paper belongs in \textit{Cognitive Science} rather than a single-discipline journal because its contribution is integrative. It brings together four traditions that have developed largely apart: conversation analysis on turn-taking, projection and collaborative completion; psycholinguistic work on interactive alignment; Clark's contribution model of grounding; and the sociolinguistic literature on conversational dominance and withdrawal. Each of these describes a feature of successful dialogue. The manuscript proposes the mechanism that explains why those features are observed together when an exchange counts as conversation, and why their absence makes the same surface activity count as something else. The predictions call for converging methods, from corpus coding and confederate paradigms to self-report and hyperscanning, which is also why the paper is aimed at a cognitive-science readership.

The manuscript is a theoretical article and reports no new data, so no ethical approval was required. I am the sole author and approve the submission. A large language model was used in preparing the manuscript, as described in the manuscript and on the title page. The manuscript has not been published and is not under consideration elsewhere.

\closing{Sincerely,}
\end{letter}
\end{document}
